\documentclass[../main.tex]{subfiles}

\begin{document}

\chapter{拓扑空间: 一般性质}\label{chapter:ch11-拓扑空间}

前两章我们研究了度量空间. 在这些空间中, 我们首先用度量定义开球, 再用开球定义开集. 我们发现可以仅用与度量相联系的开集表达一些概念. 本章我们研究基于开集的空间: 其他概念可通过开集定义. 这样的空间称为拓扑空间. 它们比度量空间更为一般. 或许你会问: 为什么不停留在度量空间? 从分析的观点来看, 主要的理由是常常需要在一个比由度量空间所提供的更为一般的框架下研究序列的收敛性或集合的紧性这样的概念. 直接的例子是考虑集合上的实值函数族. 函数序列一致收敛的概念是一个度量概念. 逐点收敛的概念不是一个度量概念. 另一个突出的例子来自作为赋范线性空间的集合 $X$. 具有由范数诱导的集合 $X$ 是一个度量空间. 关于该度量有序列收敛和集合紧性的概念. 但在 $X$ 上有一些重要的概念, 如序列的弱收敛性(我们在第 8 章研究过它)和集合的弱紧性, 它们不能在度量的框架下阐述. 它们可被阐述为拓扑赋范线性空间上的一个称为弱拓扑的拓扑概念. 此外, 不同拓扑的比较便于我们理解不同模式的序列收敛的微妙之处.

\subfile{Chapter-11/section-01.tex}

\subfile{Chapter-11/section-02.tex}

\subfile{Chapter-11/section-03.tex}

\subfile{Chapter-11/section-04.tex}

\subfile{Chapter-11/section-05.tex}

\subfile{Chapter-11/section-06.tex}

\end{document}